% =====================================================================
\section{Proposed Architecture}
\label{sec:arch}

Fig.~\ref{fig:arch} separates the architecture into the operational
path, the terminal environment, and the offline learning path. The
operational path reviews each newly announced job once and changes its
assignment through independent proposal and acceptance policies and a
central commitment procedure. The terminal environment then measures
truck dwell, additional travel, rehandling, and vessel idle cost. The
learning path clones the state and random stream immediately before a
decision and compares the observed action with counterfactual
alternatives. This computation remains outside the operational path.

\begin{center}
\begin{minipage}{\textwidth}
\centering
\includegraphics[width=\textwidth]{final_figure/system-architecture.png}
\captionof{figure}{Detailed flow of the proposed architecture. (a) Every
60\,s, the operational path checks eligibility, proposal, acceptance,
and commitment; failure of mutual consent or capacity keeps the original
assignment. (b) The discrete-event environment links the gate, 21 blocks,
two yard cranes per block, and vessel operations to calculate cost $\Phi$.
(c) Offline learning compares up to three worlds from the same state and
random stream over 3 h and constructs separate labels for the proposal
and acceptance cost networks.}
\label{fig:arch}
\end{minipage}
\end{center}

\subsection{Notation, Actions, and Policies}
\label{sec:arch-notation}
The decision interval is $\delta=60$\,s. The current assignment of job
$i$ is $z_i=(b_i,\tau_i)$, where $b_i$ is the block and $\tau_i$ is the
announced arrival time. A reallocation action is
$a_i=(b'_i,\Delta_i)\in\mathcal A_i(t)$ and gives
$z_i(a_i)=(b'_i,\tau_i+\Delta_i)$. Keeping the original assignment is
$a_i^0=(b_i,0)$. Block reallocation has $b'_i\ne b_i$, and arrival-time
adjustment has $\Delta_i>0$; a candidate may change both. Outbound jobs
are excluded from block reallocation because the retrieved container has
a fixed location.
Let $n_i$ be the announcement time, $g_i$ the gate-in time, and $r_i$ an
indicator of prior review. A job becomes eligible not when it is
announced but when its arrival comes within a review window
$w_i=\min(W,\ \tau_i-n_i)$ of the present, with $W=1800$\,s, so that the
decision is taken on the freshest state the lead time allows:

\begin{equation}
\mathcal E_t=\{i\mid w_i-\delta<\tau_i-t\le w_i,\ g_i>t,\ r_i=0\}.
\end{equation}
Each job therefore enters exactly one decision interval and is never
revisited. For a job announced at least $W$ ahead of its arrival the
decision falls exactly $W$ before that arrival---about 46 minutes later
than the median announcement, which is where the freshness comes
from---while a job announced with a shorter lead is reviewed as soon as
it is announced. A job announced no earlier than its own arrival has
$w_i=0$ and is never eligible; such jobs are 14.1\% of the demand under
the lead-time distribution of Sect.~\ref{sec:env-demand}, and the
architecture leaves their assignment untouched. The proposal
policy uses the current block state and job-action features $o_i^P$. It
scores each candidate and selects the one with the lowest predicted cost
score (Sect.~\ref{sec:arch-models}):

\begin{equation}
\hat a_i=\arg\min_{a\in\mathcal A_i(t)}
\hat c_{\theta}^{P}(o_i^P,a),\qquad
p_i=\mathbf 1[\hat a_i\ne a_i^0].
\end{equation}
If the policy proposes a change, the destination block or target time
slot becomes the accepting party. The acceptance policy uses the
proposal message and destination state $o_i^A$ to compare acceptance
with rejection:

\begin{equation}
q_i=\mathbf 1\!\left[
\hat c_{\psi}^{A}(o_i^A,\hat a_i,\mathrm{accept})
\le
\hat c_{\psi}^{A}(o_i^A,\hat a_i,\mathrm{reject})
\right].
\end{equation}
If the target resource has no remaining capacity, the request is rejected
before the network is called. The mutually accepted candidates form
$\mathcal J_t=\{i\in\mathcal E_t\mid p_iq_i=1\}$.

\subsection{Central Commitment Constraints}
\label{sec:arch-commit}
The central procedure does not create new candidates. It considers only
$\mathcal J_t$, sorted by the acceptance policy's predicted cost, with
lexicographic tie-breaking. Let $x_i\in\{0,1\}$ indicate commitment,
$F_b(t)$ denote the free slots in block $b$, and $K_k(t)$ denote the
remaining quota of time slot $k$. The procedure enforces

\begin{equation}
\begin{aligned}
x_i&\le p_iq_i,\qquad
\sum_{i\in\mathcal J_t}\!x_i\mathbf1[b'_i\!=\!b]\le F_b(t),\\
\sum_{i\in\mathcal J_t}\!x_i\mathbf1[k(i)\!=\!k]&\le K_k(t),\qquad
\sum_{t\in\mathcal T}\!\mathbf1[i\in\mathcal E_t]\le1.
\end{aligned}
\end{equation}
Block capacity is enforced in every experiment. When time-slot quota
data are unavailable, we set $K_k(t)=\infty$. Results under this setting
represent the potential range of arrival-time adjustment. All committed
assignments are applied together at the end of the decision epoch. An
earlier commitment therefore cannot change the state observed by another
candidate in the same epoch. The structure guarantees two properties:
only mutually accepted candidates are committed, and each job is
reviewed once.

\subsection{Policy Models and Inputs}
\label{sec:arch-models}
The functions $\hat c^{P}_{\theta}$ and $\hat c^{A}_{\psi}$ in
Sect.~\ref{sec:arch-notation} are cost networks. Neither produces action
probabilities. Each maps an observed state and one candidate to a scalar
cost score. The score is pair-centred rather than absolute: the two
actions of a counterfactual pair are centred on their mean before
training (Sect.~\ref{sec:arch-labels}), so the network predicts what a
candidate costs relative to the action it was rolled out against. Mean
subtraction removes a constant that is common to both actions, so the
ranking of candidates, and therefore the $\arg\min$ in~(1) and the
comparison in~(2), is unaffected. The policy evaluates every candidate
with the same network and selects the $\arg\min$ of the learned score. Both networks are multilayer perceptrons with weights
shared across blocks. The proposal network has a 21-dimensional input,
and the acceptance network has a 16-dimensional input. Both use two
64-unit hidden layers with ReLU. They contain 5{,}633 and 5{,}313
parameters, respectively, or fewer than 11{,}000 in total. Inputs
include current waiting, inventory, remaining crane work, the public
expected arrival time, route difference, and announced volume near the
target time. Unrealised arrival and completion times and simultaneous
responses from other accepting parties are excluded.
Fig.~\ref{fig:mlp} shows one cost network. A candidate enters as a
fixed-length vector $x\in\mathbb R^{d_{in}}$. It passes through three
linear layers, with ReLU after the first two, and produces one scalar:
the pair-centred cost score of that candidate. Because the output is
one value per candidate, the candidate-list length may change across
decision epochs without changing the model. The same weights also apply
to every block. The model therefore learns a common rule for evaluating
a candidate in the current block state, rather than a separate rule for
each block. The policies differ only in input dimension: 21 for proposal
and 16 for acceptance.

\begin{center}
\begin{minipage}{\textwidth}
\centering
\includegraphics[width=\textwidth]{final_figure/MLP-relu.png}

\vspace{3pt}
$\operatorname{MLP}(x)=W_3\,\operatorname{ReLU}\bigl(W_2\,\operatorname{ReLU}(W_1x+b_1)+b_2\bigr)+b_3$
\vspace{2pt}
\captionof{figure}{Structure of one cost network. The input is a fixed-length
feature vector for one candidate, and ReLU follows each of the two
64-unit hidden layers. The output is one pair-centred cost score. A
decision evaluates every candidate with the same network and takes the
$\arg\min$. The proposal and acceptance policies share this structure
and differ only in input dimension, 21 and 16 respectively.}
\label{fig:mlp}
\end{minipage}
\end{center}

\subsection{Counterfactual Cost Labels}
\label{sec:arch-labels}
For each sampled decision, we clone the state and random stream just
before branching. We then run up to three worlds for three hours:

\begin{equation}
\begin{aligned}
\phi_i^F&=\Phi_H(\text{observed}),\qquad
\phi_i^P=\Phi_H(\text{proposal reversed}),\\
\phi_i^A&=\Phi_H(\text{acceptance reversed}),
\end{aligned}
\end{equation}
with $H=10{,}800$\,s. The alternative world always reverses what the
policy actually did. If the proposal policy proposed a change, the
alternative keeps the current assignment; if it kept the assignment, the
alternative commits the cheapest candidate under the current weights.
The acceptance alternative reverses acceptance and rejection. A decision
without an acceptance request produces only the proposal comparison.
For each policy, we centre the two actions of the pair on their mean,
$\bar\phi_i^P=(\phi_i^F+\phi_i^P)/2$, and divide by KRW 100{,}000. This
gives $y_{i,F}^P=(\phi_i^F-\bar\phi_i^P)/10^5$ and
$y_{i,P}^P=(\phi_i^P-\bar\phi_i^P)/10^5$, that is, $\pm$ half the
pairwise cost difference. We apply the same centring to
$(\phi_i^F,\phi_i^A)$ for the acceptance policy. The reference point is
therefore the pair itself rather than a fixed action: the target states
how much better or worse one action was than the single alternative that
was rolled out against it in the same world. Mean subtraction removes
the cost the two worlds share without changing the order of the two
actions.

\subsection{Training Settings and Operational Separation}
\label{sec:arch-train}
We sample at most 64 decisions per day. Exploration decreases linearly
from 0.50 to 0.05, and decisions are split 80/20 into training and
validation sets. We optimise with Adam ($3\times10^{-4}$) and Huber loss
($\beta=1.0$). A few counterfactual differences are much larger than the
rest, so squared error would allow these samples to dominate the fit
~\cite{refhuber}. Minibatches contain at most 256 rows. The number of
updates is $\max(50,\min(600,4\times\text{rows}))$, and gradients are
clipped at 1.0. Counterfactual branches run in parallel across CPU cores,
and each day's labels are replaced after the daily update. During
evaluation, weights are frozen, exploration is disabled, and no
counterfactual branch is called.
The labels are therefore not a fixed dataset. Every counterfactual pair
branches from the action the current policy actually took, so the states
that are labelled, and the alternative each is compared against, both
move as the weights change; the labels of one pass are discarded after
that pass. In this respect the scheme is on-policy Monte-Carlo advantage
estimation with a simulator-generated baseline: the target is a complete
simulated rollout to the horizon, and the method uses neither
bootstrapping nor a policy gradient.
Training both networks for 30 passes took about 18 hours of wall-clock time on the machine used. Online
inference requires one forward pass through a small network for each
candidate. This cost is small relative to the accompanying constraint
check. The main computational expense is counterfactual label
generation, which is performed offline and never enters the operating
path.
